% from main.tex
%%%%%%%%%%  scan idx 609  |  folio 596  %%%%%%%%%%
\section*{1. Énoncé du théorème principal et applications}

\textbf{Théorème 1.1.} Soient \(S\) un schéma noethérien, \(X\) et \(Y\) deux \(S\)-schémas propres, \(f:X\to Y\) un morphisme. On suppose \(S\) intègre et \(f\) surjectif. Alors il existe un ouvert non-vide \(V\) de \(S\) tel que le morphisme
\[
  f^*:\Pic_{Y/S}|_V\longrightarrow \Pic_{X/S}|_V
\]
entre les foncteurs de Picard (localisés pour la topologie f.p.p.f) soit représentable par un morphisme quasi-affine de présentation finie.

\textbf{Corollaire 1.2.} Soit \(X\) un schéma propre sur un schéma noethérien \(S\). On suppose \(S\) intègre. Alors il existe un ouvert non vide \(V\) de \(S\) tel que le foncteur de Picard \(\Pic_{X/S}|_V\) soit représentable par un \(V\)-schéma en groupes somme d'ouverts quasi-projectifs de type fini, donc a fortiori séparé sur \(V\).

En effet rappelons le fait suivant (FGA N\textsuperscript{o}232)~: soient \(S\) un schéma, \(X\) un \(S\)-schéma projectif, de présentation finie et intègre sur \(S\) (c'est-à-dire plat à fibres géométriquement intègres). Alors le foncteur de Picard est représenté par un \(S\)-schéma en groupes \(\Pic_{X/S}\), réunion croissante d'ouverts quasi-projectifs de présentation finie. La considération des «~polynômes de Hilbert~» permet ici de décomposer \(\Pic_{X/S}\) en somme d'ouverts \(\Pic^q_{X/S}\), \(q\in\mathbb Q[t]\). Dans l'exposé suivant on démontre sous l'hypothèse supplémentaire \(S\) noethérien que les \(\Pic^q_{X/S}\) sont de type fini sur \(S\). On sait alors que 1.2 est vrai dans le cas où \(X\) est en plus projectif et intègre sur \(S\); on va en déduire le cas général en utilisant le

\textbf{Lemme 1.3.} Soit \(X\) de type fini sur \(S\) noethérien intègre. Alors il existe un ouvert non vide \(V\) de \(S\), un schéma intègre \(S'\) et un morphisme fini et
